\begin{tikzpicture}
    \umlactor[x=-8,y=0]{Joueur}
    \begin{umlsystem}[x=0,y=0]{Jeu Love Letter}
        \umlusecase[x=0,y=3]{Lancer le jeu}
        \umlusecase[x=0,y=1]{Quitter le jeu}
        \umlusecase[x=0,y=-1]{Jouer son tour}
        \umlusecase[x=0,y=-3]{Naviguer dans les menus}

        \umlusecase[x=10,y=3]{Consulter sa main / Score}
        \umlusecase[x=10,y=1]{Choisir une carte à jouer}
        \umlusecase[x=10,y=-1]{Piocher une carte}
        \umlusecase[x=10,y=-3]{Démarrer une partie}
    \end{umlsystem}

    \umlinclude{usecase-3}{usecase-5}
    \umlinclude{usecase-3}{usecase-6}
    \umlinclude{usecase-3}{usecase-7}
    \umlinclude{usecase-4}{usecase-8}

    \umlinherit{Joueur}{usecase-1}
    \umlinherit{Joueur}{usecase-2}
    \umlinherit{Joueur}{usecase-3}
    \umlinherit{Joueur}{usecase-4}

    \umlnote[x=0,y=5]{usecase-1}{Via Main et LoveLetter}
    \umlnote[x=16,y=1]{usecase-6}{Méthode Player.chooseCard()}
    \umlnote[x=16,y=-1]{usecase-7}{Méthode Player.drawCard()}
    \umlnote[x=0,y=-5]{usecase-4}{Via SceneManager}
\end{tikzpicture}